\subsection{The same local term and temperature can give opposite
answers}\label{proof:local_interactions:the-same-local-term-and-temperature-can-give-opposite-answers}

Compare the following at the same fixed \(\beta =6\), with the same
local term \(K(x_i,x_j)=-x_ix_j\) and positive edge weights:

\begin{center}
\begin{tabularx}{\linewidth}{@{}Xccc@{}}
\toprule
Family & Edge weights & \shortstack{Maximum absolute\\row sum} & Influence AA \\
\midrule
Disjoint matching & 1 & 1 & Holds \\
The 16-regular expanders above & \(1/16\) & 1 & Fails \\
\bottomrule
\end{tabularx}
\end{center}

For a matching, the Gibbs law is a product of two-spin blocks. Each pair
has atom probabilities \(\exp (\pm \beta )/(4\cosh  \beta )\), all at
least

\[
a_\beta=\frac1{2(e^{2\beta}+1)}>0.
\]

The bounded-block theorem in Appendix~\ref{proof:positive_structure} proves
influence AA for arbitrary independent blocks of uniformly bounded size
and an atom lower bound. It therefore applies for every fixed finite
\(\beta\), including 6. Isolated spins, if needed, are uniform singleton
blocks.

The row normalization is identical, although the edge weights differ to
compensate for degree. Both families use positive weights of the same
symmetric two-body term. The comparison is meaningful in the
positive-weight interaction setting of the supplied Marwaha--Sud paper.

This disproves an \textbf{instance-family classification that predicts a
common AA outcome solely from the local interaction term and \(\beta\)}.
It does not invalidate a worst-case classification over all allowed
interaction graphs: a worst-case family can contain easy instances and
still fail a uniform guarantee. In fact, the expander construction
proves that such a universal guarantee fails for the ferromagnetic term
at \(\beta =6\), and the gauge construction proves the same for the
antiferromagnetic term on a restricted bipartite graph class.

Both ferromagnetic families have trivially computable ground states
(all-plus and all-minus), yet they have opposite AA behavior. Thus easy
ground-state optimization does not imply a Gibbs-input influence
guarantee.

\textbf{Corollary L4 (unrestricted positive weights give a trivial
worst-case temperature boundary).} Suppose a fixed local term is either
\(-x_ix_j\) or \(+x_ix_j\), allowed with arbitrary positive edge
weights, with no uniform interaction normalization such as (N). For
every fixed \(\beta >0\), its class of Gibbs input families contains a
bounded-degree AA counterexample. At \(\beta =0\) the law is uniform and
AA holds.

To prove the positive-temperature claim, put weight \(6/(16\beta )\) on
the edges of the bipartite 16-regular expanders. For the ferromagnet,
\(\beta\) times each edge weight is exactly \(6/16\), so the Gibbs law
is the \(\beta =6\) normalized counterexample from Section~\ref{proof:local_interactions:bounded-degree-bounded-row-norm-and-an-explicit-parameter-choice}. For the
antiferromagnet use the staggered gauge. These weights are constants
with respect to n, hence meet an allowance of positive weights bounded
by \(\operatorname{poly}(n)\) for all sufficiently large n.~The
term-wise \textbf{worst-case} boundary under unrestricted scaling is
therefore zero, even though many instance families at \(\beta >0\)
satisfy AA. A meaningful nonzero temperature classification must fix a
scale or a uniform strength budget.
